%% ma: L12-B07-0032
%% lop: 12
%% chuong: 2
%% bai: 7
%% dang: 12.7.7
%% loai: TN4
%% muc: TH
%% dap_an: "A"
%% nguon: "CĐ2.6.pdf giáo viên gửi 27/10/2026 (M07, Lê Minh Khôi – Nguyễn Anh Khôi, Chương II Chủ đề 6 Ứng dụng vào các bài toán thực tế), Câu 12"
%% kien_thuc: "Bài 7 (tọa độ theo các hướng Đông, Tây, Nam, Bắc), Bài 8 (trung điểm)"
%% trang_thai: chinh-thuc
\begin{ex}
Xét hệ trục tọa độ $Oxyz$ gắn với đài kiểm soát không lưu tại $O$, mặt phẳng $(Oxy)$ là mặt đất nằm ngang, tia $Ox$ hướng Nam, tia $Oy$ hướng Đông, tia $Oz$ hướng thẳng đứng lên trời (đơn vị km). Khinh khí cầu thứ nhất nằm cách điểm xuất phát $2{,}5$ km về phía Nam, $2$ km về phía Đông và cao $0{,}8$ km. Khinh khí cầu thứ hai cách điểm xuất phát $1{,}5$ km về phía Bắc, $3$ km về phía Tây và cao $0{,}6$ km. Tọa độ trung điểm $I$ của đoạn thẳng nối hai khinh khí cầu đối với hệ trục tọa độ đã chọn là
\choice{\True $I(0{,}5;-0{,}5;0{,}7)$}{$I(2{,}0;-0{,}5;0{,}7)$}{$I(0{,}5;-1{,}0;0{,}7)$}{$I(-0{,}5;0{,}5;0{,}7)$}
\loigiai{Hai khinh khí cầu ở $(2{,}5;2;0{,}8)$ và $(-1{,}5;-3;0{,}6)$, trung điểm $I(0{,}5;-0{,}5;0{,}7)$.}
\end{ex}
